\documentclass[11pt]{article}

\usepackage[utf8]{inputenc}
\usepackage[margin=1.05in]{geometry}
\usepackage{amsmath,amssymb,amsthm,mathtools}
\usepackage{booktabs,array}
\usepackage{xcolor}
\usepackage{listings}
\usepackage[colorlinks=true,linkcolor=blue!45!black,citecolor=blue!45!black,urlcolor=blue!45!black,pdfauthor={Claude (AI model by Anthropic); submitted by Aleksandr Bersenev}]{hyperref}
\usepackage[nameinlink,capitalise,noabbrev]{cleveref}
\numberwithin{equation}{section}
\newtheorem{theorem}{Theorem}[section]
\newtheorem{proposition}[theorem]{Proposition}
\newtheorem{lemma}[theorem]{Lemma}
\newtheorem{corollary}[theorem]{Corollary}
\theoremstyle{definition}
\newtheorem{definition}[theorem]{Definition}
\theoremstyle{remark}
\newtheorem{remark}[theorem]{Remark}
\newcommand{\F}{\mathbb F}
\lstset{basicstyle=\ttfamily\small,columns=fullflexible,keepspaces=true,breaklines=true,frame=single,framerule=0.3pt,rulecolor=\color{black!30},extendedchars=true,inputencoding=utf8,
  literate={→}{{$\to$}}1 {≤}{{$\le$}}1 {≠}{{$\ne$}}1 {∀}{{$\forall$}}1 {∃}{{$\exists$}}1 {∧}{{$\wedge$}}1 {∈}{{$\in$}}1 {ℕ}{{$\mathbb N$}}1 {β}{{$\beta$}}1 {α}{{$\alpha$}}1}

\title{An improved upper bound for binary covering codes\\ of length 20 and covering radius 3}
\author{Claude (AI model by Anthropic)\\ \small run and submitted by Aleksandr Bersenev}
\date{October 14, 2026\\[3pt]\footnotesize Produced autonomously by Claude (Opus 5.5) agents in Claude Code; run and submitted by Aleksandr Bersenev.\\ \footnotesize Not an official Anthropic publication. Not peer reviewed.}

\begin{document}
\maketitle

\begin{abstract}
Let $K(n,R)$ denote the minimum number of words in a binary code of length $n$ with covering
radius $R$. We show $K(20,3)\le 1984$, improving the bound $K(20,3)\le 2048$ recorded in Kéri's
standard tables of covering codes. The code is the amalgamated direct sum of the $[7,4]$ Hamming code
with a new \emph{normal} binary code of length $14$, covering radius $2$ and $248$ words (a union of
$62$ cosets of a $[14,2,3]$ linear code). The covering property is certified in Lean~4 with Mathlib,
using only the standard axioms. We also report, as negative results, extensive but unsuccessful computer
searches aimed at the football pool problem $K_3(n,1)$ for $n=6,7,8$.
\end{abstract}

\section{Introduction}
A binary code $C\subseteq\F_2^n$ has \emph{covering radius} $R$ if every word of $\F_2^n$ is within
Hamming distance $R$ of some codeword. The quantity $K(n,R)$, the minimum size of such a code, has been
studied since the 1960s; see the monograph of Cohen, Honkala, Litsyn and Lobstein~\cite{CHLL} and the
football pool problem $K_3(n,1)$ for its ternary analogue~\cite{vLAvLW89,LMT09}. The best known bounds
for $n\le 33$, $R\le 10$ are collected in G.~Kéri's tables~\cite{Keri}, which have been the standard
reference since the 2000s (the binary table we consulted is dated October~2009; its update log ends in
November~2011). In these tables
\[
  892\;\le\;K(20,3)\;\le\;2048,
\]
where the upper bound carries no reference mark, i.e.\ it is classified as ``trivial'' (inherited from
neighbouring entries).

\paragraph{Result.} We prove $K(20,3)\le1984$ (\cref{thm:main}). The construction is classical --- the
amalgamated direct sum (ADS) of Graham and Sloane~\cite{GS85} and Cohen, Lobstein and Sloane~\cite{CLS86}
--- and the observation behind it is simple: the ADS of the $[7,4]$ Hamming code with a \emph{normal}
code of length $14$ and covering radius $2$ having $M$ words has length $20$, covering radius $3$ and
$8M$ words. With $M=256=2\cdot128$ (twice the tabulated bound for $K(13,2)$) this gives exactly the recorded value
$2048$, which suggests how the entry may have arisen; the table entry $K(14,2)\le248$ (credited in~\cite{Keri} to Bertolo, Di Pasquale and Santisi, 2006) would give $1984$,
provided some $248$-word code is normal --- a property that is not automatic and that, as far as we can
tell, was never established or used. We construct such a normal code directly
(\cref{sec:B}) and verify the whole argument in Lean.

To our knowledge, as of October 14, 2026, no upper bound below $2048$ for $K(20,3)$ has been published;
we searched the update log of~\cite{Keri} and recent work on covering codes
\cite{GijswijtPolak25,Marosi26,DMP25}, which concerns lower bounds, alphabets $q\ge5$ and linear
codes respectively. A systematic check of all binary and mixed binary/ternary table entries under direct
sums, puncturing, extension and ADS (assuming normality) found no other entry that should have improved
in this way.

This is a modest improvement of one table entry. The work started as an attempt at the football pool
problem (\cref{sec:negative}); we report those negative results as well.

\section{Definitions}
For $x,y\in\F_2^n$ let $d(x,y)$ be the Hamming distance and $d(x,C)=\min_{c\in C}d(x,c)$.
For a code $C$ and $s\in\{0,1\}$ write $C_i^s=\{c\in C: c_i=s\}$.

\begin{definition}[Normality]
A code $C$ of covering radius $R$ is \emph{normal with respect to coordinate $i$} if
$d(x,C_i^0)+d(x,C_i^1)\le 2R+1$ for all $x\in\F_2^n$.
\end{definition}

\begin{definition}[Amalgamated direct sum]
Let $A\subseteq\F_2^{n_A}$ and $B\subseteq\F_2^{n_B}$. The ADS (amalgamated on the last coordinate of
$A$ and the first coordinate of $B$) is
\[
  A\dot\oplus B=\{(a_1,\dots,a_{n_A-1},s,b_2,\dots,b_{n_B}) :
  (a_1,\dots,a_{n_A-1},s)\in A,\ (s,b_2,\dots,b_{n_B})\in B\}\subseteq\F_2^{n_A+n_B-1}.
\]
\end{definition}
Clearly $|A\dot\oplus B|=|A^0||B^0|+|A^1||B^1|$, where $A^s$ (resp.\ $B^s$) is taken with respect to the
amalgamated coordinate.

\begin{proposition}[\cite{GS85,CLS86}]\label{prop:ads}
If $A$ has covering radius $R_A$ and is normal w.r.t.\ its last coordinate, and $B$ has covering radius
$R_B$ and is normal w.r.t.\ its first coordinate, then $A\dot\oplus B$ has covering radius at most
$R_A+R_B$.
\end{proposition}
\begin{proof}
Write $x=(\alpha,t,\beta)$ with $\alpha\in\F_2^{n_A-1}$, $t\in\F_2$, $\beta\in\F_2^{n_B-1}$, and
for $s\in\{0,1\}$ choose $a^s\in A^s$, $b^s\in B^s$ nearest to $(\alpha,t)$ and $(t,\beta)$. The word
$c^s$ obtained by amalgamating $a^s$ and $b^s$ lies in $A\dot\oplus B$ and
$d(x,c^s)=d((\alpha,t),a^s)+d((t,\beta),b^s)-[t\neq s]$. Summing over $s$ and using normality,
$d(x,c^0)+d(x,c^1)\le(2R_A+1)+(2R_B+1)-1$, so $\min_s d(x,c^s)\le R_A+R_B$.
\end{proof}

\section{A normal code of length 14 and covering radius 2 with 248 words}\label{sec:B}
Number the coordinates $0,\dots,13$ and identify a word with the integer $\sum_i x_i2^i$.
Let $H=[I_{12}\mid h_{12}\ h_{13}]$ be the $12\times14$ binary matrix whose columns are
$e_0,\dots,e_{11}$, $h_{12}=e_1+e_2$ and $h_{13}=e_0+e_1+e_2$ (as integers: $1,2,4,\dots,2048,6,7$).
Its kernel is the $[14,2,3]$ code spanned by $e_1+e_2+e_{12}$ and $e_0+e_1+e_2+e_{13}$. Let $S\subseteq\F_2^{12}$
be the following set of $62$ syndromes (integers, bit $j$ = row $j$):
\begin{center}\small
\begin{minipage}{0.92\linewidth}\raggedright
4, 74, 188, 231, 322, 377, 401, 428, 553, 599, 655, 738, 783, 820, 986, 993, 1055, 1138, 1194, 1236,
1313, 1468, 1487, 1562, 1601, 1719, 1785, 1879, 1900, 1922, 2097, 2114, 2220, 2271, 2330, 2343, 2505,
2546, 2575, 2684, 2706, 2791, 2897, 2922, 2948, 3007, 3183, 3201, 3258, 3338, 3346, 3420, 3556, 3575,
3620, 3665, 3788, 3903, 3905, 3993, 4015, 4082.
\end{minipage}
\end{center}
Define $B=\{x\in\F_2^{14}: Hx\in S\}$, a union of $62$ cosets of $\ker H$, so $|B|=248$.

\begin{proposition}\label{prop:B}
$B$ has covering radius $2$, $|B^0|=|B^1|=124$ with respect to coordinate $0$, and $B$ is normal with
respect to coordinate $0$: $d(y,B^0)+d(y,B^1)\le5$ for all $y\in\F_2^{14}$.
\end{proposition}
The proof is a finite computation over the $2^{14}$ words (it is part of the Lean certificate,
\cref{sec:lean}; independently, a breadth-first computation of $d(\cdot,B^0)$ and $d(\cdot,B^1)$ gives
$\max_y(d(y,B^0)+d(y,B^1))=5$). In fact $B$ is normal with respect to every coordinate.
The size $248$ equals the best known upper bound for $K(14,2)$~\cite{Keri}; it does not improve it.

\paragraph{How $B$ was found.} Codes of the form $\{x:Hx\in S\}$ with $H=[I_k\mid h_{k}\ \cdots\ h_{n-1}]$
have covering radius $\le R$ iff $S$ is an $R$-dominating set in the Cayley graph of $\F_2^k$ generated by
the columns of $H$ (the ``matrix method''). For $n=14$, $k=12$ there are only $239$
classes of the two extra columns up to coordinate permutations (given by the numbers of coordinates in
which only $h_{12}$, only $h_{13}$, or both are $1$). We ran a weighted local search for small
$2$-dominating sets on each class; the class $(0,1,2)$, i.e.\ $h_{12}=e_1+e_2$, $h_{13}=e_0+e_1+e_2$,
gave $|S|=62$ within seconds (it is the only class reaching $62$ in our runs; $194$ classes reach $64$), and normality was then checked.

\section{The main theorem}
Let $A$ be the $[7,4]$ Hamming code with parity-check columns $1,2,\dots,7$ (coordinate $j$ has column
$j+1$):
$A=\{0,7,25,30,42,45,51,52,75,76,82,85,97,102,120,127\}$. It is perfect, $|A^0|=|A^1|=8$ with respect to
coordinate $6$, and normal with respect to every coordinate (one checks
$d(y,A^0)+d(y,A^1)\le3$ for all $y\in\F_2^7$).

\begin{theorem}\label{thm:main}
$C=A\dot\oplus B$, amalgamated on coordinate $6$ of $A$ and coordinate $0$ of $B$, is a binary code of
length $20$ with $1984$ words and covering radius $3$. Hence $K(20,3)\le1984$.
\end{theorem}
\begin{proof}
Length $7+14-1=20$; size $8\cdot124+8\cdot124=1984$; covering radius $\le1+2$ by
\cref{prop:ads,prop:B}. (A word at distance exactly $3$ exists, e.g.\ by the sphere bound.)
In integer form, $C=\{(u \bmod 64)+64v : u\in A,\ v\in B,\ \lfloor u/64\rfloor = v\bmod 2\}$.
\end{proof}

\begin{remark}
Equivalently, $C=\{x\in\F_2^{20}: H_Cx\in\{0\}\times S\}$, where $H_C$ is the $15\times20$ matrix whose first
three rows are the parity checks of $A$ on coordinates $0,\dots,6$ and whose last twelve rows are $H$ on
coordinates $6,\dots,19$; thus $C$ is a union of $62$ cosets of a $[20,5]$ linear code.
\end{remark}

\begin{remark}
A brute-force check over all $2^{20}$ words confirms that the $1984$ words are distinct and that every
word is within distance $3$ of the code. The data files accompany this note (\texttt{data/}).
\end{remark}

\section{Formal verification}\label{sec:lean}
The file \texttt{lean/Main.lean} (Lean 4.35.0, Mathlib) proves
\begin{lstlisting}
theorem main : ∃ C : Finset (Fin 20 → Fin 2), C.card ≤ 1984 ∧
    ∀ x : Fin 20 → Fin 2, ∃ c ∈ C, hammingDist x c ≤ 3
\end{lstlisting}
and the companion statement \texttt{codeB} for $B$ (length $14$, $\le248$ words, covering radius $\le2$).
Here \texttt{hammingDist} is Mathlib's Hamming distance (number of coordinates where two functions
differ), so the statement is literally $K(20,3)\le1984$.

\paragraph{Structure of the proof.} Words are encoded by natural numbers $a<2^n$ through their binary
digits (\texttt{dec n a}); the lemma \texttt{hamming\_dec} proves that \texttt{hammingDist} of two
decoded words equals a recursively defined digit-difference count \texttt{dist n a b}, and
\texttt{dec\_surj} (via Mathlib's \texttt{finFunctionFinEquiv}) shows that every word is decoded from
some $a<2^n$. The finite facts --- normality of $A$ over $128$ words and normality of $B$ over $2^{14}$
words, the latter with packed witness pairs checked by a tail-recursive Boolean function whose soundness
lemma is proved --- are established by \texttt{decide +kernel}, i.e.\ evaluated by the Lean kernel; no
\texttt{native\_decide} is used. The ADS covering argument of \cref{prop:ads} is proved in Lean
(lemmas \texttt{key} and \texttt{cover\_nat}).

\paragraph{Axioms.} \texttt{\#print axioms} reports \texttt{[propext, Classical.choice, Quot.sound]} for
both \texttt{main} and \texttt{codeB}. Checking took $105$--$151$\,s wall time and about $12$\,GB peak memory on one
core.

\paragraph{Statement audit.} \texttt{Fin 20} $\to$ \texttt{Fin 2} is the set of binary words of length $20$;
\texttt{C.card} $\le$ \texttt{1984} together with the covering condition is exactly the definition of
$K(20,3)\le1984$. The code is not hidden behind a definition in the statement: the existential is
discharged by an explicit list.

\section{Negative results on the football pool problem}\label{sec:negative}
Our primary target was an upper bound in the football pool problem $K_3(n,1)$, whose records
$K_3(6,1)\le73$, $K_3(7,1)\le186$, $K_3(8,1)\le486$ date from 1989~\cite{vLAvLW89}. We did not improve any
of them. For the record:
\begin{itemize}
\item $n=7$: the matrix method with $k=6$ (all six inequivalent choices of the seventh column) yields at
best $3\cdot62=186$; with $|S|=61$ our local search always left at least one vertex uncovered.
The methods with $k=5,4$ gave at best $189$ and $216$. We also treated every conjugacy class of
fixed-point-free automorphisms of order $3$ of the Hamming graph $H(7,3)$ (products of $c\in\{0,1,2\}$
coordinate $3$-cycles with a symbol translation on $w\ge1$ fixed coordinates; $12$ classes): codes
invariant under such an automorphism correspond to dominating sets of a $729$-vertex quotient graph, and
the best size found was again $3\cdot62=186$. Allowing a few non-invariant extra words (``almost
invariant'' codes: $60$ or $61$ orbits plus an exactly optimised set of extra words covering the
remaining uncovered words) never went below $186$. A weighted local search on the full space $\F_3^7$
reached only $192$--$193$; codes invariant under random subgroups of the automorphism group (orbit local
search and integer programming) also failed.
\item $n=8$: the matrix method reproduces $486$ ($k=6,|S|=54$ and $k=4,|S|=6$) but nothing smaller.
\item $n=6$: $73$-word codes are found in milliseconds, $72$ never.
\end{itemize}
A sweep over the open entries of Kéri's binary, ternary, quaternary and mixed tables with full-space local
search reproduced table values for small parameters and never improved one; the records reflect
algebraic structure that unstructured search does not find.

\section{Open questions}
Any normal binary code of length $14$, covering radius $2$ and $M<248$ words would give
$K(20,3)\le 8M$ and at the same time improve $K(14,2)\le248$. Our searches (all two-extra-column
classes of the matrix method with $k=12$, one-extra-column classes with $k=13$, and codes formed by $123$
cosets of a $[14,1]$ code plus a few extra words) produced nothing below $248$; the best unions of cosets
of a one-dimensional code have $250$ words. Similarly, $K(19,3)\le1024$ and $K(21,3)\le3072$ coincide
with $8K(13,2)$ and $8K(15,2)$, so improved normal codes for $K(13,2)$ or $K(15,2)$ would improve those
entries as well. The density of our code, $1984\cdot\sum_{i\le3}\binom{20}{i}/2^{20}\approx2.56$, is
still larger than that of the neighbouring entries ($\approx2.27$ for $n=19$ and $\approx2.29$ for
$n=21$), which suggests that $K(20,3)$ can be improved further by a direct construction.

\section{Limitations}
The improvement is small and rests on a classical construction. Novelty is asserted only to our
knowledge: Kéri's tables were last updated around 2009--2011 and an unpublished or obscure later
improvement of $K(20,3)$ cannot be completely excluded. The Lean proof certifies the theorem
independently of the search code.

\begin{thebibliography}{99}
\bibitem{CHLL} G.~Cohen, I.~Honkala, S.~Litsyn, A.~Lobstein, \emph{Covering Codes}, North-Holland, 1997.
\bibitem{CLS86} G.~D.~Cohen, A.~C.~Lobstein, N.~J.~A.~Sloane, Further results on the covering radius of
codes, \emph{IEEE Trans.\ Inform.\ Theory} 32 (1986) 680--694.
\bibitem{GS85} R.~L.~Graham, N.~J.~A.~Sloane, On the covering radius of codes, \emph{IEEE Trans.\ Inform.\
Theory} 31 (1985) 385--401.
\bibitem{Keri} G.~Kéri, Tables for bounds on covering codes,
\url{https://old.sztaki.hu/~keri/codes/index.htm} (binary table dated 2009-10-15; accessed 2026-10-14).
\bibitem{vLAvLW89} P.~J.~M.~van Laarhoven, E.~H.~L.~Aarts, J.~H.~van Lint, L.~T.~Wille, New upper bounds for
the football pool problem for 6, 7, and 8 matches, \emph{J.\ Combin.\ Theory Ser.\ A} 52 (1989) 304--312.
\bibitem{LMT09} J.~Linderoth, F.~Margot, G.~Thain, Improving bounds on the football pool problem by
integer programming and high-throughput computing, \emph{INFORMS J.\ Comput.} 21 (2009) 445--457.
\bibitem{GijswijtPolak25} D.~Gijswijt, S.~Polak, Semidefinite lower bounds for covering codes,
arXiv:2504.01932 (2025, rev.\ 2026).
\bibitem{Marosi26} M.~Marosi, New upper and lower bounds on covering codes $K_q(n,R)$ for alphabets of
size $5\le q\le21$, arXiv:2608.19872 (2026).
\bibitem{DMP25} A.~A.~Davydov, S.~Marcugini, F.~Pambianco, New upper bounds for binary linear covering
codes, arXiv:2511.02542 (2025).
\end{thebibliography}

\end{document}
